\documentclass[main]{subfiles}
\usepackage[sols]{workbook}
\numbering{ws2-3}

\begin{document}

\ifSubfilesClassLoaded{%
  \chaptermark{\chaptertwo}
}{}

\section{Taylor's Remainder Theorem} \label{ws2-3}

\begin{framed}
\centerline{\textbf{Lesson Goals}}

You should be able to:
\begin{itemize}

\item Use the remainder theorem to find an upper bound for the error in approximating a function by a degree $n$ Taylor polynomial
\item Use the remainder theorem to find the degree of a Taylor polynomial that leads to a desired margin of error
\end{itemize}
\end{framed}



 \begin{problemonly} 
    \begin{framed}
    \textbf{Reminder:} The $n$th degree Taylor polynomial for $f(x)$ centered at $c$ is given by
    $$
      P_{n}(x) = f(c) + f'(c)(x-c) + \frac{f''(c)}{2!}(x-c)^{2} + \cdots + \frac{f^{(n)}(c)}{n!}(x-c)^{n}.
    $$
    \end{framed}
  \end{problemonly}

\begin{enumerate}
   
 \item \label{motivate-taylor}
    You are asked to predict the time dependence of sales of the iPhone X. The theoretical models from
economics give a complicated function $f(t)=\sin(t)+t^2$ with $t$ given in months.
\begin{enumerate}
    \item \begin{problem}[\vfill]
        Find the second order Taylor expansion for $f$ centered at 0.
    \end{problem}
    
    \begin{solution}
        $P_2(t)=t+t^2$.
    \end{solution}
    \item \begin{problem}[\vfill]
        Here are approximate values of the $\sin$ function at various $t$: $\sin(\frac{1}{2})\approx .48,\sin(1) \approx 0.84$.
Use these values to estimate the error when using the second order Taylor approximation at $t=\frac{1}{2}$ and $t=1$. By this, I mean that you should compute the value of the second order Taylor
approximations at $t =\frac{1}{2}$
and $t = 1$ and see directly how much they differ from the values given above for the $\sin$ function at these $t$-values.

    \end{problem}
    
    \begin{solution}
        The error is $|f(t)-P_2(t)|$.
        $f(\frac{1}{2})\approx 0.73$ and $P_2(\frac{1}{2})=.75$, so the error is approximately $0.02$.

        $f(1)\approx 1.84$ and $P_2(1)=2$ so the error is approximately $0.16$.
    \end{solution}
    \item \begin{problem}[\vfill]
        We expect the size of the error to be on the order of the next term in the Taylor approximation. Compute this next term and compare its size with the error.
    \end{problem}

    \begin{solution}
        The next term in Taylor's expansion is $\frac{1}{6}f^{(3)}(0)t^3=\frac{-1}{6}t^3$. At $t=\frac{1}{2}$ this has an absolute value of approximately $0.021$. At $t=1$, it has absolute value approximately $0.167$.
    \end{solution}

\end{enumerate}
\pspace{\newpage}
  \begin{problemonly} 
  \begin{framed}
  \textbf{Remainder Theorem.} Let $f$ be a function whose $n$-th Taylor polynomial is $P_n(x)=\displaystyle \sum_{k = 0}^{n}
    \frac{f^{k}(c)}{k!}(x-c)^{k}$.
    and suppose $M_{(n+1)}$ is an upper bound for $|f^{(n+1)}(x)|$ on the interval between $x$ and $c$.  Then
    $$
      |f(x) - P_{n}(x)| \leq \frac{M_{(n+1)}}{(n+1)!}|x-c|^{n+1} 
    $$
\end{framed}
\end{problemonly}
\item \begin{enumerate}
    \item \begin{problem}[\vfill]
        Use the first order Taylor approximation based at $x = 100$ to estimate $\sqrt{102}$ and $\sqrt{99}$.
    \end{problem}

    \begin{solution}
        $P_1(x)=10+\frac{1}{20}(x-100)$.
        Then $P_1(102)=10.1$ and $P_1(99)=9.95$.
    \end{solution}
\item \begin{problem}[\vfill]
    We shall now find an upper bound for the difference between your estimates in (a) and the true
value of these square roots. First, write the error bound from Taylor’s theorem in our case.
\end{problem}

\begin{solution}
    The error bound in our case is $\frac{M_{(2)}|x-100|^2}{2}$.
\end{solution}
\item \begin{problem}[\vfill]
    Now find a viable choice for $M_{(2)}$ in the case $t = 102$. Then, use this value of $M_{(2)}$ with $t = 102$ in the error bound formula.
\end{problem}

\begin{solution}
    The second derivative is $f''(x)=\dfrac{-1}{4x^{\frac{3}{2}}}$ which is largest at $x=100$, which makes $\frac{1}{4000}$ a reasonable choice for $M_{(2)}$. Plugging this into our error formula, $\text{error}\leq \dfrac{2^2}{2\cdot 4000}=\dfrac{1}{2000}$.
\end{solution}
\item\begin{problem}[\vfill]
    Find a viable choice for $M_{(2)}$ in the case $x = 99$. Then use this value of $M_{(2)}$ and also $x = 99$ in the error bound formula.
\end{problem} 
\begin{solution}
    The second derivative is $f''(x)=\dfrac{-1}{4x^{\frac{3}{2}}}$ which is largest at $x=99$. However, we don't know $\sqrt{99}$ (this is what we are approximating), so we can choose something nearby, like $81$ which makes $\frac{1}{4\cdot 729}$ a reasonable choice for $M_{(2)}$. To simplify computations further, we could say this is less than $\frac{1}{2500}$. Using this value for $M_{(2)}$, we see our error$\leq \dfrac{1}{2\cdot2500}=\dfrac{1}{5000}$.
\end{solution}
\end{enumerate}

\pspace{\newpage}

\item\begin{problem}[\vfill]
For $t$ near 0, the first order Taylor approximation based at $t = 0$ for sin(t) is $P_1(t) = t$. What is a
reasonable choice for $M_{(2)}$ in this case if we want to bound the error between $\sin(t)$ and $t$ for $0\leq t \leq 1$?    
\end{problem} 
\begin{solution}
    The second derivative of $\sin t$ is $-\sin t$ so a reasonable choice for $M_{(2)}$ is 1.
\end{solution}

\item \begin{enumerate}
\item \begin{problem}[1in]
    Write out the Taylor series for $\cos(x)$.
\end{problem}

\begin{solution}
    The Taylor series for $\cos x$ is $\displaystyle\sum_{k=0}^\infty \dfrac{(-1)^kx^{2k}}{(2k)!}$
\end{solution}
\item \begin{problem}[\vfill]
    What is $M_{(n+1)}?$
\end{problem}

\begin{solution}
    We know any derivative for $\cos x$ is $\pm \sin x$ or $\pm \cos x$. All of these functions are bounded by 1, so 1 is a reasonable choice for $M_{(n+1)}$.
\end{solution}
\item\begin{problem}[\vfill]
    We are approximating $\cos(\frac{1}{4})$. Give a reasonable $n$ so that the error is definitely less than $2^{-20} $ (which is roughly $10^{-6}$).
\end{problem} 

\begin{solution}
    We see that $|f(x)-P_n(x)|\leq \dfrac{(\frac{1}{4})^{n+1}}{(n+1)!}$. We want this expression to be less than $2^{-20}$ to ensure that the actual error is also smaller than that. Notice that $4^{n+1}=2^{2n+2}$. So we are looking for an $n$ such that $\dfrac{1}{(n+1)!2^{2n+2}}\leq \dfrac{1}{2^{20}}$. If we ignore the $(n+1)!$ (because this only makes the left side smaller, we could find an $n$ such that $\dfrac{1}{2^{2n+2}}\leq \dfrac{1}{2^{20}}$. We simply solve $2n+2=20$, and see that $n=9$ works. Because of the $(n+1)!$ a smaller choice of $n$ also works, indeed $n=6$ is the smallest value for $n$, but 9 is a reasonable choice.
\end{solution}
\end{enumerate}
  
\pspace{\newpage}

  \item Let $f(x) = \sin(2x)$.  
  \begin{enumerate}
    \item
    \begin{problem}[2.0in]
      Write down the Taylor series for $f$ centered at $0$, using sigma notation.  
    \end{problem}


    \begin{solution}
 Here are the first five derivatives of $f$:
    $$
      f'(x) = 2 \cos(2x) \qquad f''(x) = -4 \sin(2x) \qquad f'''(x) = -8 \cos(2x) $$
      
      $$f^{(4)}(x) = 16 \sin(2x) \qquad f^{(5)}(x) = 2^{5} \cos(2x)  $$
    If you evaluate each of these at $0$, notice that $f^{(n)}(0) = 0$ whenever $n$ is even.
    The Taylor series for $f$ centered at $0$ is:
    $$
      \frac{2}{1!} (x-0)^{1} - \frac{8}{3!} (x-0)^{3} + \frac{2^{5}}{5!} (x-0)^{5} \cdots
    $$
    In sigma notation, the Taylor series for $f(x) = \sin(2x)$ is
    $$
      \sum_{k = 0}^{\infty} (-1)^{k} \frac{2^{2k+1}}{(2k+1)!} x^{2k + 1}.
    $$   
    \end{solution}

    \item
    \begin{problem}
      Let $P_n$ denote the degree $n$ Taylor polynomial for $f$, centered at $0$.
      Using the Remainder Theorem,
      find a reasonable choice of $n$ for which $|P_{n}(x) - f(x)| \leq 10^{-3}$ for every $x$
      in the interval $\displaystyle \left[ - \frac{1}{2}, \: \frac{1}{2} \right]$.
     
    \end{problem}  
      
    \begin{solution}
      We need a reasonable choice for $M_{(n+1)}$, an upper bound for
      $|f^{(n+1)}(x)|$ on the interval $[- \frac{1}{2}, \: \frac{1}{2} ]$.  Since
      the sine and cosine functions are bounded between $-1$ and $1$, from our above
      computations of derivatives we see that $M_{(n+1)} = 2^{n+1}$ is a safe choice.
      The Remainder Theorem then tells us that
      $$ 
        |f(x) - P_{n}(x)| \leq \frac{2^{n+1}}{(n+1)!} \: |x - 0|^{n+1}
      $$     
      for every $x$ in the interval $[- \frac{1}{2}, \: \frac{1}{2} ]$.  For all
      values of $x$ in this interval, notice that $\displaystyle |x - 0|^{n+1} \leq \frac{1}{2^{n+1}}$.      
      Thus, our error bound tells us that
      $$ 
        |f(x) - P_{n}(x)| \leq \frac{2^{n+1}}{(n+1)!} \: \frac{1}{2^{n+1}} \; = \; \frac{1}{(n+1)!}
      $$
      for all $x$ in the interval $[- \frac{1}{2}, \: \frac{1}{2}]$.  If we want the assurance
      that $|f(x) - P_{n}(x)| \leq 10^{-3}$, we can guarantee this by insisting that
      $$
        \frac{1}{(n+1)!} \leq \frac{1}{10^3}.  
      $$
      Equivalently, we need $(n+1)! \geq 1000$.  We don't need the very smallest choice of
      $n$ that satisfies this (which happens to be $n = 6$).  For example, if we use $n = 11$
      then $(n+1)! = 12! = (12)(11)(10) 9!$ is definitely larger than 1000 because $(12)(11)(10) > 10^3$.
    \end{solution}  
  \end{enumerate}



 \end{enumerate}



\end{document}